\documentclass[a4paper,11pt]{article}
\usepackage[utf8]{inputenc}
\usepackage[T1]{fontenc}
\usepackage[english]{babel}
\usepackage[left=2.5cm,right=2.5cm,top=2cm,bottom=2cm]{geometry}
\usepackage{hyperref}
\usepackage{xcolor}
\usepackage{microtype}
\usepackage{lmodern}

\definecolor{primary}{RGB}{0, 85, 164}
\hypersetup{colorlinks=true, urlcolor=primary}

\begin{document}
\noindent
\textbf{Lo\"ic Aron MBASSI EWOLO} \\
ENSEA -- Double Dipl\^ome Ing\'enieur \\
\href{mailto:loicaron.mbassiewolo@ensea.fr}{loicaron.mbassiewolo@ensea.fr} \quad (+33) 7 59 52 33 98 \\
\href{https://github.com/Nameless0l}{github.com/Nameless0l} \quad \href{https://mbassiloic.tech}{mbassiloic.tech}

\vspace{1em}
\noindent
Cergy, le \today

\vspace{1em}
\noindent
\textbf{Objet : Candidature -- Data Scientist, IA \& Recommandation} \\
\textbf{R\'ef\'erence : Offre d'alternance 12 mois -- BOULANGER}

\vspace{1.5em}
\noindent Madame, Monsieur,

\vspace{0.8em}
\noindent Leader français du retail électronique en transformation numérique, Boulanger poursuit une ambition majeure en data science et IA générative pour optimiser recommandations produits, personnalisation client et expérience d'achat. Je suis particulièrement attiré par votre approche centrée sur l'intelligence de recommandations multi-critères et par votre volonté d'exploiter l'IA pour créer une expérience client exceptionnelle. C'est pourquoi je candidats avec conviction pour rejoindre votre équipe en tant que Data Scientist IA et Recommandation en alternance dès septembre 2026.

\vspace{0.8em}
\noindent J'ai conçu et déployé un système intelligent de recommandations multi-agents orches-trées par Google ADK : quatre agents collaboratifs avec LLM Gemini pour identifier besoins agriculteurs (Planning, Météo, Marché, Recommandations). Cette architecture produit des recommandations multi-critères (rendement, rentabilité, durabilité) en intégrant dynamiquement des données externes (APIs OpenWeather, marchés agricoles). Cette expérience directe de systèmes de recommandation me forme aux enjeux concrets : pertinence contextuelle, agrégation multi-critères, ranking algorithmique et engagement utilisateur. Je peux transposer cette architecture à vos contextes retail/e-commerce Boulanger pour recommandations produits personnalisées et pertinentes.

\vspace{0.8em}
\noindent Depuis avril 2026, je déploie un pipeline NLP complet à l'équipe TRIBE (INRIA Saclay) : fine-tuning BERT production, classification multi-labels, scoring interprétable (0-100). Cette compétence NLP directement applicable à l'analyse de reviews clients, sentiment analysis et enrichissement de descriptions produits pour Boulanger. Ma recherche chez MILA (Montréal) sur la généralisation extrême de réseaux neuronaux m'a formé aux fondations deep learning et aux limites de modèles, essentielles pour valider la robustesse avant déploiement. En outre, mes succès en compétitions data (1er prix IndabaX Africa en IA santé, 1er prix CITS national en optimisation) démontrent ma capacité à délivrer solutions data science complexes et impactantes, du prototypage au déploiement production.

\vspace{0.8em}
\noindent Je suis disponible dès septembre 2026 pour une alternance 12 mois, ce qui m'permettra de délivrer des initiatives data science concrètes pour Boulanger (systèmes recommandation, analyse client, optimisation conversion). Mes expériences en recommandations multi-critères, NLP et optimisation algorithmique me positionnent idéalement pour accélérer votre transformation data science. Je serais ravi de discuter comment mon expertise en systèmes recommandation intelligents, mes compétences NLP et ma vision d'une IA personnalisante correspondent à l'ambition d'excellence client de Boulanger. Merci pour cette opportunité exceptionnelle.

\vspace{1.5em}
\noindent Veuillez agr\'eer, Madame, Monsieur, l'expression de mes salutations distingu\'ees.

\vspace{2em}
\noindent \textbf{Lo\"ic Aron MBASSI EWOLO}
\end{document}
